\section{Scriptural loci and their theological use}
\label{app:scripture}

The table groups the principal passages developed in the body. It is an index
of this study's exposition, not an inventory of every biblical appearance of
an angel. Psalm numbers follow the Vulgate, with the Hebrew number in brackets.

\begin{longtable}{@{}>{\raggedright\arraybackslash}p{0.31\linewidth}>{\raggedright\arraybackslash}p{0.63\linewidth}@{}}
\toprule
Locus & Subject in the study\\
\midrule\endhead
Genesis 3:24 & Cherubim and the guarded way to life; the names of the orders.\\
Genesis 18:1--15; 19:1--3 & Divine visitation, hospitality and assumed bodies.\\
Genesis 28:10--17 & Jacob's ladder; angels and God's covenant promise.\\
Exodus 23:20--23 & The angel's mission and the authority of the sender.\\
Tobias 12:12--22 & Raphael, intercession, healing and the return to praise of God.\\
Psalms 90:11--12 [91]; 102:20--22 [103] & Guardianship, obedience and service.\\
Psalm 148:1--5 & Heavenly praise and the angels' created existence.\\
Isaias 6:1--8 & Seraphim, the burning coal, purification and mission;
Dionysius's difficult case in chapter XIII.\\
Daniel 7:9--10 & The innumerable court; distinction between standing before
God and ministering.\\
Daniel 10:12--21; 12:1--3 & Michael, the princes of peoples, providential
government and human resurrection.\\
Matthew 18:10 & The little ones and their angels; guardianship and vision
of the Father.\\
Matthew 22:29--32 & Risen humanity compared to angels concerning marriage;
human resurrection remains human.\\
Luke 1:19--38; 2:8--15 & Gabriel, the Annunciation and the heavenly praise
at Christ's birth.\\
Luke 24:1--8 & The messengers at the tomb and the announcement of the Resurrection.\\
Colossians 1:15--20; 2:18--19 & Invisible powers created through Christ;
the primacy of the head over every ministry.\\
Hebrews 1:1--14 & The Son adored by angels; ministering spirits.\\
Hebrews 12:22--24 & The heavenly assembly of angels and perfected human
spirits around the mediator.\\
Apocalypse 12:7--12 & Michael's battle, the accuser's defeat and the Lamb's victory.\\
Apocalypse 19:9--10; 22:8--9 & The angel as fellow servant; adoration directed to God.\\
\bottomrule
\end{longtable}

\section{Demonic action and its limits in Aquinas}
\label{app:powers}

These distinctions summarize the determinations expounded in the treatise.
They concern the kinds of action Aquinas attributes to fallen spirits,
not the identification of a cause in any particular human experience.

\begin{longtable}{@{}>{\raggedright\arraybackslash}p{0.22\linewidth}>{\raggedright\arraybackslash}p{0.49\linewidth}>{\raggedright\arraybackslash}p{0.22\linewidth}@{}}
\toprule
Subject & Determination and limit & Governing locus\\
\midrule\endhead
Natural intelligence & Natural cognition remains; charity and the blessed
vision do not follow from it. & I q.~64 a.~1\\
The will & The fallen will is fixed in its evil end; its natural powers do
not thereby become infinite. & I q.~64 a.~2\\
Suffering & Sorrow concerns a will thwarted in its objects, not the passions
of a mortal body. & I q.~64 a.~3\\
Order & Natural superiority and subordination can remain among demons;
their concord in evil is not charity. & I q.~109 aa.~1--2\\
Illumination & Their communication is not the holy illumination treated
under the good angels. & I q.~109 a.~3\\
Subjection & Evil spirits remain subject to the ordering exercised by good
angels under God. & I q.~109 a.~4\\
Material effects & Angels move bodies and employ bodily causes; matter does
not receive a substantial form by an angel's mere will. & I q.~110 aa.~2--3\\
Miracles & A marvel exceeding familiar human power need not exceed all
created power. A miracle strictly so called is God's work. & I q.~110 a.~4;
q.~114 a.~4\\
Human will & A spirit can persuade and stir passions; it cannot inwardly
cause or necessitate a person's consent. & I q.~111 a.~2\\
Imagination and sense & Created spiritual agency can affect bodily
conditions of imagination and sensation; this does not imply access to
every secret thought. & I q.~111 aa.~3--4;
q.~57 a.~4\\
Temptation & The devil tempts toward sin, but not every sin results from
a present demonic suggestion. & I q.~114 aa.~2--3\\
Providence & Permission of trial and the tempter's evil intention are
distinct; the good intended by God does not make the temptation good. &
I q.~114 a.~1\\
Resistance & Aquinas describes restraint following resistance; he does not
promise that a person can never be tempted again. & I q.~114 a.~5\\
\bottomrule
\end{longtable}

The theological account of evil is accompanied by the Church's confession
that the devil remains a creature (\emph{Catechism}, 395). The history,
rite and discipline of exorcism receive their separate treatment in
Triptych's \emph{Catholic Exorcism: History and Current Practice}.
